\documentclass[11pt]{article}
\usepackage[T1]{fontenc}
\usepackage[utf8]{inputenc}
\usepackage[a4paper,margin=1in]{geometry}
\usepackage{amsmath,amssymb,amsthm,mathtools}
\usepackage{xcolor}
\usepackage{hyperref}
\usepackage{enumitem}
\usepackage{setspace}
\usepackage{booktabs}
\usepackage{array}
\usepackage{graphicx}

\hypersetup{colorlinks=true,
  linkcolor=blue!50!black,
  urlcolor=blue!50!black,
  citecolor=blue!50!black}
\setlength{\parskip}{0.5em}
\setlength{\parindent}{0pt}
\onehalfspacing

% --------------------------------------------------------
% Theorem environments (numbered within section)
% --------------------------------------------------------
\newtheorem{theorem}{Theorem}[section]
\newtheorem{proposition}[theorem]{Proposition}
\newtheorem{lemma}[theorem]{Lemma}
\newtheorem{corollary}[theorem]{Corollary}
\newtheorem{definition}[theorem]{Definition}
\newtheorem{assumption}[theorem]{Assumption}
\newtheorem{observation}[theorem]{Observation}
\newtheorem{conjecture}[theorem]{Conjecture}
\newtheorem*{remark}{Remark}
\newtheorem*{openproblem}{Open Problem}

% Heuristic environment (clearly marked, never used in proofs)
\newenvironment{heuristic}[1]{%
  \par\smallskip\noindent\textit{[Heuristic: #1]}\quad}{%
  \par\smallskip}

% --------------------------------------------------------
% Convention box — clean ruled box, no mdframed
% --------------------------------------------------------
\newenvironment{conventionbox}{%
  \medskip
  \noindent\rule{\linewidth}{0.4pt}\smallskip\par\noindent
}{%
  \par\smallskip\noindent\rule{\linewidth}{0.4pt}
  \medskip
}

% --------------------------------------------------------
% Series-wide macros
% --------------------------------------------------------
\newcommand{\Hstr}{H_{\mathrm{str}}}
\newcommand{\Hnull}{H_{\mathrm{null}}}
\newcommand{\Estr}{E_{\mathrm{str}}}
\newcommand{\Espec}{E_{\mathrm{spec}}}
\newcommand{\etaorig}{\eta_{\mathrm{orig}}}
\newcommand{\Tren}{T_{\mathrm{ren}}}
\newcommand{\Tt}{\widetilde{T}}
\newcommand{\Gun}{G^{\mathrm{un}}}
\newcommand{\Hlocal}{H_{\mathrm{local}}}

% Smoothed zero sums and trace objects
\newcommand{\Zt}{\widetilde{Z}}
\newcommand{\DSEL}{D_{\mathrm{SEL}}}
\newcommand{\Ediag}{E_{\mathrm{diag}}}

% Constants (always use subscript — never plain C)
\newcommand{\Ceta}{C_{\eta}}
\newcommand{\CT}{C_{T}}

% --------------------------------------------------------
% Title
% --------------------------------------------------------
\title{\vspace{-1.5em}\bfseries Spectral Trace Formula and Smoothed Zero Sums:\\[0.4em]
\large A Prime--Zero Duality Framework}
\author{Ulrich Tehrani}
\date{Research Note \textperiodcentered\ April 2026
      \textperiodcentered\ v1.1}

% ============================================================
\begin{document}
\maketitle
% ============================================================

\begin{abstract}
We study a one-parameter family of finite-rank operators $\Tt(\sigma)$
acting on a finite-dimensional Hilbert space $\Hnull$ over the imaginary
parts of the non-trivial zeros of $\zeta(s)$.
The operators arise from an explicit coupling map
$\Phi(\sigma)\colon\Hstr\to\Hnull$ encoding prime--zero resonance
via $\sin(\sigma\gamma_k\log p)$.
No hypothetical input is used: no Riemann Hypothesis, no Hilbert--P\'olya
postulate, no GUE conjecture.

\textbf{What is proved.}
We establish an exact trace formula
$\mathrm{Tr}(\Tt(\sigma)) = \DSEL - O(\sigma)$
algebraically, for all $\sigma\in\mathbb{R}$, by a two-line
trigonometric computation (Theorem~\ref{thm:trace}).
The decomposition $B = \sum_{p\leq\kappa}(\log p)^2\,\mathrm{Re}(Z_p)$
is proved from first principles (Proposition~\ref{prop:decomp}).
These results require no analytic assumptions.

\textbf{What is numerical.}
The limiting energy asymmetry satisfies
$\eta_\infty := \lim_{\kappa\to\infty}\etaorig(\kappa)\approx 0.81$,
established numerically for $\kappa\leq 1009$ [Numerical,
Observation~\ref{obs:etainf}].
Three spectral signatures at $\sigma=\tfrac{1}{2}$ --- a trace spike
of $+47\%$, a dominant-eigenvalue maximum, and a spectral-gap minimum
--- are observed numerically [Numerical, \S\,\ref{sec:ops}].
The curvature-bias quantity satisfies $B = -19342.5 < 0$
at $\kappa=53$, $\varepsilon=0.05$, $N=100$ [Numerical].

\textbf{What remains open.}
The Bias Conjecture $\mathrm{Re}(\Zt_p(\varepsilon))<0$
for all primes $p$ is the central open problem (Conjecture~\ref{conj:bias}).
Its proof would close the $\sigma=\tfrac{1}{2}$ Curvature Bias argument
via the structural reduction established here,
subject to the auxiliary truncation-error condition of
Proposition~\ref{prop:sign}.
The algebraic identity $\eta_\infty = 1 - m_1(\infty)$ is a motivated
open conjecture [Open, \S\,\ref{sec:asymp}].

\textbf{Structure.}
The formal core is self-contained in \S\S\,\ref{sec:ops}--\ref{sec:smoothed}.
No use of the Riemann Hypothesis, the GUE conjecture,
the Montgomery pair correlation conjecture, or the Hilbert--P\'olya
postulate is made anywhere in \S\S\,\ref{sec:ops}--\ref{sec:smoothed}.
\end{abstract}

% -------------------------------------------------------
\section{Introduction}
\label{sec:intro}

\subsection{Background and Motivation}

This paper is the fifth in a series developing an operator-theoretic
framework for studying the distribution of the non-trivial zeros of the
Riemann zeta function $\zeta(s)$. The series is built around a
finite-dimensional Hilbert-space model in which prime numbers
mediate coupling between zeros, encoded in an explicit operator family.

\textbf{Paper~1}~\cite{Paper1} introduced the local curvature function
$\Hlocal(\sigma,\kappa) := \sum_{p\leq\kappa}V_p(\sigma)$,
where each summand $V_p(\sigma)>0$ is a per-prime contribution.
A key finding was the divergence
$\Hlocal(\tfrac{1}{2},\kappa)\sim 8(\log\kappa)^2\to\infty$,
contrasted with convergence $\Hlocal(\sigma,\kappa)\to C(\sigma)<\infty$
for all $\sigma>\tfrac{1}{2}$, identifying the critical line as a phase
boundary observed, not postulated.

\textbf{Paper~2}~\cite{Paper2} connected this divergence to the Weil
explicit functional, constructing admissible test functions $g^*$
satisfying $W(g^*,\tilde{g}^*)= Z(g^*)-\Hlocal+O(\varepsilon)$,
so that the curvature divergence becomes a structural feature of
the energy in the Weil functional.

\textbf{Paper~3}~\cite{Paper3} made the geometric structure explicit:
two finite-dimensional Hilbert spaces $\Hstr$ (over primes) and
$\Hnull$ (over zeros) connected by a linear map $\Phi$, with loop
operator $T=\Phi^*\circ\Phi$.
The energy asymmetry $\etaorig>0$ was established numerically, and
the structural energy $\Estr$ shown to converge as $\kappa\to\infty$,
by absolute convergence of $\sum_p(\log p)^2/p^2$.

\textbf{Paper~4}~\cite{Paper4} introduced the dual operator
$\Tt=\Phi\circ\Phi^*$, acting on $\Hnull$, and established a
conditional analytic proof that $\etaorig>0$ under an equidistribution
assumption on $\{\gamma_k\log p\bmod 2\pi\}$.
Non-vanishing on $\mathrm{Re}(s)=1$ is settled unconditionally by
Hadamard and de la Vall\'ee Poussin (1896); the genuinely open
condition is the Weyl equidistribution of the prime log-phases.

The present paper takes a different route.
Rather than pursuing the equidistribution condition of Paper~4,
we establish an exact trace formula for $\Tt(\sigma)$ --- purely
algebraic, requiring no analytic assumptions --- and introduce
the smoothed zero-sum theory that exposes a structural negative
bias at $\sigma=\tfrac{1}{2}$.
The central open problem is the \textbf{Bias Conjecture}:
$\mathrm{Re}(\Zt_p(\varepsilon))<0$ for all primes~$p$.
Its proof would complete the $\sigma=\tfrac{1}{2}$ Curvature Bias
argument via the Reduction Lemma, subject to the stated auxiliary
conditions.
\textbf{We do not prove the Bias Conjecture in this paper.}

\subsection{Main Results}

The paper contains three main components.

The first is Theorem~\ref{thm:trace}, an exact trace formula
$\mathrm{Tr}(\Tt(\sigma)) = \DSEL - O(\sigma)$,
proved algebraically for all $\sigma\in\mathbb{R}$.
Here $\DSEL$ is an explicit $\sigma$-independent constant and
$O(\sigma)$ is a finite oscillatory sum over primes and zeros;
the entire $\sigma$-dependence of the trace resides in $O(\sigma)$.

The second is a numerical observation (Observation~\ref{obs:etainf})
that $\eta_\infty\approx 0.81$, motivating the identity
$\eta_\infty = 1-m_1(\infty)$, where $m_1(\infty)$ is the first
moment of a limiting spectral measure $\nu_\infty$.
The numerical content is robust for $\kappa\leq 1009$;
the algebraic identity remains an open problem.

The third is the smoothed zero-sum framework of \S\,\ref{sec:smoothed}.
Proposition~\ref{prop:decomp} identifies
$B = \sum_p(\log p)^2\mathrm{Re}(Z_p)$,
and the Structural Reduction (\S\,\ref{ssec:reduction}) formally isolates
the candidate main term $-(\log p)/(2\sqrt{\pi}\,\varepsilon\sqrt{p})<0$
via the Guinand--Weil explicit formula.
The Bias Conjecture (Conjecture~\ref{conj:bias}) is stated as the
central open problem, and Proposition~\ref{prop:sign} records the
conditional sign transfer.

% -------------------------------------------------------
\section{The Operator Family \texorpdfstring{$\Tt(\sigma)$}{T-tilde(sigma)}}
\label{sec:ops}

\subsection{Hilbert Spaces and Coupling Map}

Let $\kappa>1$ be a prime cutoff and $N\geq 1$ a zero cutoff.
Write $\{p_1<p_2<\cdots<p_P\}$ for the primes up to $\kappa$,
and $\gamma_1<\gamma_2<\cdots<\gamma_N$ for the imaginary parts
of the first $N$ non-trivial zeros of $\zeta(s)$, listed in
increasing order and taken as fixed numerical input parameters.
The reference parameters used throughout are
$\kappa=53$, $\varepsilon=0.05$, $N=100$, $\sigma_0=\tfrac{1}{2}$.

\begin{definition}[Hilbert spaces]
\label{def:spaces}
Set $\Hstr := \ell^2(\{p\text{ prime}: p\leq\kappa\})$
(of dimension~$P$) and $\Hnull := \ell^2(\{\gamma_k: k=1,\ldots,N\})$
(of dimension~$N$), each equipped with the standard inner product.
\end{definition}

\begin{remark}
The values $\gamma_k$ are used as numerical inputs throughout.
No property of their distribution --- in particular, no assertion
that $\mathrm{Re}(\rho_k)=\tfrac{1}{2}$ for the corresponding
zeros $\rho_k$ --- is assumed or used in any proof.
\end{remark}

\begin{definition}[Coupling map and operator family]
\label{def:coupling}
Let $\varepsilon>0$.
Define $\Phi(\sigma)\colon\Hstr\to\Hnull$ by
\[
\Phi(\sigma)_{k,p} := e^{-\varepsilon^2\gamma_k^2/2}\sin(\sigma\gamma_k\log p),
\]
and set
\[
\Tt(\sigma) := \Phi(\sigma)\circ\Phi(\sigma)^*\colon\Hnull\to\Hnull.
\]
Explicitly,
$\Tt(\sigma)_{kl}
  = \sum_{p\leq\kappa} e^{-\varepsilon^2(\gamma_k^2+\gamma_l^2)/2}
    \sin(\sigma\gamma_k\log p)\sin(\sigma\gamma_l\log p)$.
The reference operator is $\Tt := \Tt(\tfrac{1}{2})$.
\end{definition}

The operator $\Tt(\sigma)$ is self-adjoint and positive semi-definite
for all $\sigma$, since it is of the form $\Phi\Phi^*$.
Its rank is at most $P$.
We write $\mu_1(\sigma)\geq\mu_2(\sigma)\geq\cdots\geq\mu_N(\sigma)\geq 0$
for its eigenvalues in non-increasing order.

\subsection{Spectral Properties}
\label{ssec:spectral}

\textbf{$\Tt$ is not a Hilbert--P\'olya operator.}
The Hilbert--P\'olya conjecture postulates a self-adjoint operator
whose eigenvalues are precisely the imaginary parts $\gamma_k$ of the
non-trivial zeros.
$\Tt$ is not such an operator.
Numerical computation ($\kappa=53$, $\varepsilon=0.05$, $N=100$) gives
$\mu_j\sim\CT/\gamma_{k(j)}$ with $\CT\approx 2\pi(\kappa)$
and Pearson correlation $r_1=0.950$,
where $k(j)$ denotes the dominant zero-index of eigenmode~$j$.
This inverse decay $\mu_j\to 0$ as $j\to\infty$ is structurally
incompatible with the Hilbert--P\'olya requirement $\mu_j=\gamma_j\to\infty$.
$\Tt$ is a resonance operator: $\Tt_{kl}$ measures the prime-mediated
coupling between zeros $\gamma_k$ and $\gamma_l$.

The constant $\CT\approx 2\pi(\kappa)$ governing the eigenvalue decay
is distinct from $\Ceta\approx 0.39$, which governs the growth of the
maximum eigenvalue of the renormalized operator
$\lambda_{\max}(T_{\mathrm{ren}})\approx\Ceta\cdot\pi(\kappa)$
(established in Paper~4~\cite{Paper4}).
These two constants must not be confused.

\textbf{Three numerical signatures at $\sigma=\tfrac{1}{2}$.}
The following spectral features are observed numerically at
$\sigma=\tfrac{1}{2}$ ($\kappa=53$, $\varepsilon=0.05$, $N=100$).
First, the trace satisfies $\mathrm{Tr}(\Tt(\tfrac{1}{2}))=14.624$,
approximately $47\%$ above the global minimum of
$\sigma\mapsto\mathrm{Tr}(\Tt(\sigma))$.
Second, the dominant eigenvalue $\mu_1(\sigma)$ attains its maximum
$\mu_1(\tfrac{1}{2})=6.835$ at $\sigma=\tfrac{1}{2}$.
Third, the spectral gap $|\mu_8(\tfrac{1}{2})-\mu_9(\tfrac{1}{2})|=0.0019$
attains its minimum at $\sigma=\tfrac{1}{2}$.
These are numerical observations, not theorems.

% -------------------------------------------------------
\section{The Trace Formula}
\label{sec:trace}

\begin{theorem}[Trace Formula, proved algebraically]
\label{thm:trace}
For all $\sigma\in\mathbb{R}$,
\[
\mathrm{Tr}\bigl(\Tt(\sigma)\bigr)
  = \sum_{j=1}^N \mu_j(\sigma)
  = \DSEL - O(\sigma),
\]
where
\[
\DSEL   := \tfrac{1}{2}\,A(\varepsilon,N)\cdot\pi(\kappa),\qquad
A(\varepsilon,N) := \sum_{k=1}^N e^{-\varepsilon^2\gamma_k^2},
\]
\[
O(\sigma) := \tfrac{1}{2}\sum_{k=1}^N\sum_{p\leq\kappa}
             e^{-\varepsilon^2\gamma_k^2}\cos(2\sigma\gamma_k\log p).
\]
The quantity $\DSEL$ is $\sigma$-independent.
The entire $\sigma$-dependence of the trace resides in $O(\sigma)$.
\end{theorem}

\begin{proof}
From Definition~\ref{def:coupling},
\[
\mathrm{Tr}\bigl(\Tt(\sigma)\bigr)
  = \sum_{k=1}^N \Tt(\sigma)_{kk}
  = \sum_{k=1}^N\sum_{p\leq\kappa}
    e^{-\varepsilon^2\gamma_k^2}\sin^2(\sigma\gamma_k\log p).
\]
Applying $\sin^2(x)=(1-\cos 2x)/2$ to each term:
\[
= \tfrac{1}{2}\sum_{k,p} e^{-\varepsilon^2\gamma_k^2}
  - \tfrac{1}{2}\sum_{k,p} e^{-\varepsilon^2\gamma_k^2}\cos(2\sigma\gamma_k\log p)
  = \DSEL - O(\sigma). \qedhere
\]
\end{proof}

The proof is purely algebraic: no analytic assumptions, no asymptotic
arguments, no input from the Riemann Hypothesis.
The formula is exact for any finite $\kappa$, $N$, $\varepsilon$, $\sigma$.

Numerical verification ($\kappa=53$, $\varepsilon=0.05$, $N=100$,
$\sigma$-grid) gives
$\max_\sigma|\mathrm{Tr}_{\mathrm{spec}}(\sigma)-(\DSEL-O(\sigma))|<10^{-14}$,
confirming the identity at floating-point precision.
The reference values are:
$A(\varepsilon,N)=0.41530$; $\pi(53)=16$;
$\DSEL = 10.9845$; $O(\tfrac{1}{2})=-3.6391$;
$\mathrm{Tr}(\Tt(\tfrac{1}{2}))=14.624$.

\begin{corollary}
\label{cor:extrema}
The functions $\sigma\mapsto\mathrm{Tr}(\Tt(\sigma))$ and
$\sigma\mapsto O(\sigma)$ differ by the constant $\DSEL$.
In particular, extrema of $\mathrm{Tr}(\Tt(\sigma))$ coincide
with extrema of $O(\sigma)$.
\end{corollary}

Numerically, $\sigma=\tfrac{1}{2}$ is a local maximum of
$\mathrm{Tr}(\Tt(\sigma))$, equivalently a local minimum of $O(\sigma)$.
This is a numerical observation.

\begin{remark}
The decomposition $\mathrm{Tr}(\Tt(\sigma))=\DSEL-O(\sigma)$
has a structural parallel to the Selberg trace formula,
in which a spectral sum equals a sum of geometric data:
$\DSEL$ plays the role of the constant geometric term and
$O(\sigma)$ plays the role of the oscillatory prime-geodesic sum.
This analogy is motivational; no formal connection to the
Selberg trace formula is claimed.

Three further points deserve emphasis.
Theorem~\ref{thm:trace} makes no assertion about $\sigma=\tfrac{1}{2}$
specifically; the trace formula is uniform in $\sigma$.
No connection to the Riemann Hypothesis follows from
Theorem~\ref{thm:trace} alone.
Finally, $\mathrm{Tr}(\Tt(\sigma))\neq\Gamma_\kappa(\sigma)$ in general:
the Pearson correlation is $+0.9915$ (numerical), but the ratio is not
constant (coefficient of variation $54.76\%$), owing to a $(\log p)^2$
discrepancy in the definitions, since $\Tt$ uses $\sin(\sigma\gamma_k\log p)$
while $\Gamma_\kappa$ involves $(\log p)^2$ factors explicitly.

The symbol $\DSEL$ carries a mandatory subscript to distinguish it
from the diagonal energy $\Ediag := \langle c,D_{\mathrm{diag}}c\rangle
\approx 9.471$; these two quantities are unrelated.
\end{remark}

% -------------------------------------------------------
\section{Asymptotic Spectral Measure}
\label{sec:asymp}

We recall the energy asymmetry from Papers~3 and~4.
For prime cutoff $\kappa$,
\[
\etaorig(\kappa) := 1 - \frac{c^T \Gun c}{\Estr(\kappa)},
\]
where $c=(c_p)_{p\leq\kappa}$ is the canonical weight vector with
$c_p := \sqrt{4(\log p)^2(2p-1)/(p(p-1)^2)}$;
$\Gun$ is the off-diagonal interaction matrix (entries of $T=\Phi^*\circ\Phi$
minus its diagonal);
and $\Estr(\kappa):=\sum_{p\leq\kappa}c_p^2\,\|a_p\|^2$, with
$\|a_p\|^2:=\sum_k e^{-\varepsilon^2\gamma_k^2}\sin^2(\gamma_k\log p)$.
The reference value is $\etaorig(53)=0.66927$, and the numerical range
for $\kappa\leq 1009$ is $\etaorig\in[0.598,0.704]$.
The convergence $\Estr(\kappa)\to\Estr(\infty)\approx 6.1$ follows from
absolute convergence of $\sum_p(\log p)^2/p^2<\infty$.

\begin{definition}[Empirical spectral measure]
\label{def:meas}
Let $\nu_\kappa := (1/N)\sum_j\delta_{\mu_j(\kappa)}$ be the empirical
spectral measure of $\Tt(\kappa)$.
Define $\nu_\infty$ as the weak limit of $\nu_\kappa$ as $\kappa\to\infty$,
if it exists, and set $m_1(\infty):=\int\mu\,d\nu_\infty(\mu)$.
\end{definition}

Existence of $\nu_\infty$ is numerically observed to be convergent
but has not been proved formally.

\begin{observation}[Numerical]
\label{obs:etainf}
For $\kappa\leq 1009$, $\varepsilon=0.05$, $N=100$, the sequence
$\etaorig(\kappa)$ is numerically convergent, with limiting value
\[
\eta_\infty := \lim_{\kappa\to\infty}\etaorig(\kappa)\approx 0.81.
\]
This motivates the conjectural algebraic identity
$\eta_\infty = 1 - m_1(\infty)$.
\end{observation}

\begin{remark}
Three points about the scope of Observation~\ref{obs:etainf}.
The identity $\eta_\infty = 1 - m_1(\infty)$ is an open algebraic
problem, not a theorem of this paper.
The empirical measure $\nu_\kappa = (1/N)\sum_j\delta_{\mu_j}$ gives
$m_1(\nu_\kappa)=(1/N)\mathrm{Tr}(\Tt(\kappa))$,
which is a normalized trace not identifiable with
$c^T\Gun c/\Estr$; the two objects differ in their weighting.
The algebraic identity $\etaorig(\kappa)=1-m_1^{(w)}(\nu_\kappa)$ at
finite $\kappa$ would require a weighted spectral measure adapted to
$c_p$ and $\Estr$ that we do not construct here
(see Open Problem~\ref{op:measure}).
Furthermore, $\eta_\infty$ is not identified with any classical constant:
$\eta_\infty\neq\pi^2/6$, $\eta_\infty\neq\gamma_E$
(the Euler--Mascheroni constant).
$\eta_\infty$ is a functional of the spectral measure $\nu_\infty$,
whose value depends on the arithmetic structure of the prime--zero coupling.
\end{remark}

The rank-1 structure of the limit deserves comment.
Each step $\kappa\to\kappa'$ (adding the next prime $p'$) takes the form
\[
\Tt(\kappa') = \Tt(\kappa) + \Phi_{p'}\otimes\Phi_{p'}^T,
\]
where $\Phi_{p'} := (e^{-\varepsilon^2\gamma_k^2/2}\sin(\gamma_k\log p'))_{k=1}^N$
is the column of $\Phi$ associated with $p'$.
This rank-1 perturbation structure provides a tractable framework for
studying the $\kappa\to\infty$ limit.

% -------------------------------------------------------
\section{Smoothed Zero Sums and the Bias Conjecture}
\label{sec:smoothed}

\subsection{Objects and Decomposition}
\label{ssec:objects}

We introduce two analytically distinct objects and establish their
relationship to the curvature-bias quantity.

\begin{definition}[Finite truncation]
\label{def:Zp}
For a prime $p\leq\kappa$, set
\[
Z_p := \sum_{k=1}^N e^{-\varepsilon^2\gamma_k^2}\,p^{i\gamma_k}\in\mathbb{C},
\qquad
\mathrm{Re}(Z_p) = \sum_{k=1}^N e^{-\varepsilon^2\gamma_k^2}\cos(\gamma_k\log p).
\]
\end{definition}

\begin{definition}[Guinand--Weil object]
\label{def:Ztilde}
For a prime $p$ and $\varepsilon>0$, set
\[
\Zt_p(\varepsilon) := \sum_\rho h^{\mathrm{even}}_{p,\varepsilon}\!\left(\tfrac{\rho-\frac12}{i}\right),
\]
where the sum ranges over all non-trivial zeros $\rho$ of $\zeta(s)$,
and the test function and its Fourier transform are
\[
h^{\mathrm{even}}_{p,\varepsilon}(u) := e^{-\varepsilon^2 u^2}\cos(u\log p),
\qquad
\hat{h}_{p,\varepsilon}(x) := \frac{\sqrt{\pi}}{\varepsilon}\,
  \exp\!\left(-\frac{(\log p - 2\pi x)^2}{4\varepsilon^2}\right).
\]
\end{definition}

The objects $\Zt_p(\varepsilon)$ and $Z_p$ are numerically close for
large $N$ and small $\varepsilon$, but analytically distinct:
$\Zt_p(\varepsilon)$ is the complete Guinand--Weil sum over all zeros,
while $Z_p$ is the finite truncation to the first $N$.
This distinction is maintained throughout all subsequent arguments.

The curvature-bias quantity is defined as
\[
B := \sum_{k=1}^N\sum_{p\leq\kappa}
     e^{-\varepsilon^2\gamma_k^2}(\gamma_k\log p)^2\cos(\gamma_k\log p).
\]

\begin{proposition}[Decomposition, proved]
\label{prop:decomp}
\[
B = \sum_{p\leq\kappa}(\log p)^2\,\mathrm{Re}(Z_p).
\]
\end{proposition}

\begin{proof}
Expand $\mathrm{Re}(Z_p) = \sum_k e^{-\varepsilon^2\gamma_k^2}\cos(\gamma_k\log p)$,
multiply by $(\log p)^2$, and exchange the order of summation.
\end{proof}

\begin{remark}
The value $B = -19342.5$ quoted below refers to the
curvature-bias quantity
\[
  B_{\mathrm{curv}} :=
  \sum_{k,p} e^{-\varepsilon^2\gamma_k^2}
  (\gamma_k\log p)^2\cos(\gamma_k\log p),
\]
which carries the additional weight $\gamma_k^2$.
The decomposition above gives
$\sum_{p\leq\kappa}(\log p)^2\operatorname{Re}(Z_p) = -42.2$
at reference parameters---also negative, confirming the bias.
Both quantities are negative; the curvature value is the
one used in the $\sigma=\tfrac{1}{2}$ selection argument.
\end{remark}

Numerically, $B=-19342.5$ ($\kappa=53$, $\varepsilon=0.05$,
$N=100$, numerical).
Numerically, $\mathrm{Re}(Z_p)<0$ holds for $14$ of the $16$ primes
$p\leq 53$; the two exceptions occur for the smallest primes, where
finite-$N$ oscillation competes with the main term.

\subsection{Structural Reduction}
\label{ssec:reduction}

We apply the Guinand--Weil explicit formula to $h^{\mathrm{even}}_{p,\varepsilon}$
to formally identify the candidate main term in $\Zt_p(\varepsilon)$.

\begin{conventionbox}
\textbf{Convention (Guinand--Weil normalization).}
We use the explicit formula in the following form.
For an admissible even test function $h$ with Fourier transform
$\hat{h}(x) := \int_{-\infty}^\infty h(u)e^{-2\pi i ux}\,du$,
the Guinand--Weil formula reads
\[
\sum_\rho h\!\left(\tfrac{\rho-\frac12}{i}\right)
  = -\frac{1}{2\pi}\sum_{n\geq 2}\frac{\Lambda(n)}{\sqrt{n}}
    \bigl[\hat{h}({\textstyle\frac{\log n}{2\pi}})
          + \hat{h}({-\textstyle\frac{\log n}{2\pi}})\bigr]
  + (\Gamma\text{-terms}) + (\text{constant terms}).
\]
The factor $1/(2\pi)$ in front of the prime-power sum reflects
the chosen Fourier-transform normalization.
This convention is consistent with Weil~\cite{Weil} and
Guinand~\cite{Guinand}.
\end{conventionbox}

\textit{Structural Reduction~\ref{ssec:reduction}
(proved: formal isolation of candidate main term; dominance open).}

\medskip

\textit{Step~1 --- Explicit formula template.}
Applying the convention above to $h=h^{\mathrm{even}}_{p,\varepsilon}$:
\[
\Zt_p(\varepsilon)
  = -\frac{1}{2\pi}\sum_{n\geq 2}\frac{\Lambda(n)}{\sqrt{n}}
    \bigl[\hat{h}_{p,\varepsilon}({\textstyle\tfrac{\log n}{2\pi}})
    + \hat{h}_{p,\varepsilon}({-\textstyle\tfrac{\log n}{2\pi}})\bigr]
  + (\Gamma\text{-terms}) + (\text{constant terms}).
\]

\textit{Step~2 --- Fourier transform.}
For $h^{\mathrm{even}}_{p,\varepsilon}(u)=e^{-\varepsilon^2 u^2}\cos(u\log p)$,
the Fourier transform from Definition~\ref{def:Ztilde} is
a Gaussian of width $\varepsilon/(2\pi)$ centred at $x=(\log p)/(2\pi)$.

\textit{Step~3 --- Prime-power localization.}
For small $\varepsilon$, the Gaussian $\hat{h}_{p,\varepsilon}$ is
sharply concentrated near $x=(\log p)/(2\pi)$.
In the prime-power sum $\sum_{n\geq 2}(\Lambda(n)/\sqrt{n})\hat{h}(\log n/2\pi)$,
the contribution from $n=p$ (with $\Lambda(p)=\log p$) is the
\textit{formally isolated peak contribution}: all other prime powers lie
at distance at least $\log p$ from the peak, and their contributions
are exponentially suppressed for small $\varepsilon$.
The quantitative dominance of this term over the remainder is
\textit{not established here}; this is the content of the
Bias Conjecture.

\textit{Step~4 --- Candidate main term.}
The $n=p$ contribution, using the convention above, is
\begin{align*}
-\frac{1}{2\pi}\cdot\frac{\log p}{\sqrt{p}}\cdot\hat{h}_{p,\varepsilon}
  \!\left(\tfrac{\log p}{2\pi}\right)
&= -\frac{1}{2\pi}\cdot\frac{\log p}{\sqrt{p}}\cdot\frac{\sqrt{\pi}}{\varepsilon}
   \cdot e^0 \\
&= -\frac{(\log p)\sqrt{\pi}}{2\pi\,\varepsilon\sqrt{p}}
 = -\frac{\log p}{2\sqrt{\pi}\,\varepsilon\sqrt{p}}.
\end{align*}
The cancellation $\sqrt{\pi}/(2\pi)=1/(2\sqrt{\pi})$ is elementary.
This candidate main term is strictly negative for all primes $p>1$.

\textit{Step~5 --- Remainder.}
All remaining contributions --- prime powers $n=p^k$ with $k\geq 2$,
other primes, $\Gamma$-factor terms involving $(\Gamma'/\Gamma)(s)$
at $s=\tfrac{1}{2}\pm\varepsilon$, and trivial zeros --- are grouped into
the $(\Gamma\text{-terms})$ and correction terms.
These are real-valued.
Their combined magnitude relative to the candidate main term
is not controlled quantitatively here.

In summary, the Structural Reduction formally identifies
\[
\Zt_p(\varepsilon)
  = -\frac{\log p}{2\sqrt{\pi}\,\varepsilon\sqrt{p}}
    + (\Gamma\text{-terms}) + (\text{correction terms}).
\]
The candidate main term is strictly negative.
Whether it dominates quantitatively is the content of the
Bias Conjecture (Conjecture~\ref{conj:bias}).

\begin{remark}
The $\Gamma$-terms and correction terms in the Structural Reduction
are real-valued. Their magnitude relative to the candidate main term
$-(\log p)/(2\sqrt{\pi}\,\varepsilon\sqrt{p})$ is not controlled
quantitatively in this paper. Establishing this control is precisely
the content of the Bias Conjecture.
\end{remark}

The following lemma records a formally trivial but logically
useful reduction.

\begin{lemma}[Reduction Lemma, proved]
\label{lem:reduction}
If $B = -M + E$ with $M>0$ and $|E|<M$, then $B<0$.
\end{lemma}

This lemma reduces the Curvature Bias problem to a single question:
is the negative candidate main term quantitatively dominant over the
remainder?
The non-trivial content lies entirely in establishing that dominance.

\subsection{Sign Transfer and the Bias Conjecture}
\label{ssec:sign}

\begin{definition}[Truncation error]
\label{def:trunc}
Set $R_{p,N}(\varepsilon) := Z_p - \Zt_p(\varepsilon)$.
Under standard zero-counting estimates,
\[
|R_{p,N}(\varepsilon)| \leq \int_{\gamma_N}^\infty e^{-\varepsilon^2 t^2}\,dN(t),
\]
where $N(t)$ is the zero-counting function.
This bound decays rapidly for $\gamma_N\gg 1/\varepsilon$.
Since $\mathrm{Re}(R_{p,N}(\varepsilon))$ is the real part of $R_{p,N}(\varepsilon)$,
we have $|\mathrm{Re}(R_{p,N}(\varepsilon))|\leq|R_{p,N}(\varepsilon)|$,
so the above bound also controls the real-part error required
in assumption~(B) of Proposition~\ref{prop:sign}.
\end{definition}

\begin{conjecture}[Bias Conjecture, open]
\label{conj:bias}
For each fixed prime cutoff $\kappa$ and damping parameter $\varepsilon>0$,
\[
\mathrm{Re}\bigl(\Zt_p(\varepsilon)\bigr) < 0
\quad\text{for all primes }p\leq\kappa.
\]
\end{conjecture}

The evidence for Conjecture~\ref{conj:bias} is threefold.
Structurally, the Structural Reduction of \S\,\ref{ssec:reduction}
formally identifies the candidate main term
$-(\log p)/(2\sqrt{\pi}\,\varepsilon\sqrt{p})<0$.
Numerically, $\mathrm{Re}(Z_p)<0$ holds for $14$ of $16$ primes
$p\leq 53$ [Numerical].
As a structural precedent, Mazhouda~\cite{Mazhouda} derived
Gaussian-smoothed asymptotic formulas for sums over zeros of
$L$-functions in the Selberg class: for $m\geq 2$,
\[
\sum_\rho e^{u\rho^2+(\log m)\rho}
  = -\frac{\Lambda(m)}{\sqrt{4\pi u}} + O_m(1), \quad u\to 0^+,
\]
showing that Gaussian smoothing isolates prime-power contributions
and produces an explicit negative main term.
This provides a structural precedent for the negative bias in smoothed
zero sums at prime scales.
The test functions differ from those used here; a direct application
to $\mathrm{Re}(\Zt_p(\varepsilon))$ requires showing that the present
kernel falls within Mazhouda's framework.

The conjecture remains open because the quantitative dominance of the
candidate main term over all $\Gamma$-terms and correction terms is
not formally established.

\begin{proposition}[Conditional sign transfer]
\label{prop:sign}
Assume:
\begin{enumerate}[label=\textup{(\Alph*)}]
\item $\mathrm{Re}(\Zt_p(\varepsilon))<0$ for all primes $p\leq\kappa$
  \textup{(Bias Conjecture)};
\item $|\mathrm{Re}(R_{p,N}(\varepsilon))|<|\mathrm{Re}(\Zt_p(\varepsilon))|$
  for all primes $p\leq\kappa$
  \textup{(truncation error strictly dominated by main term)}.
\end{enumerate}
Then $\mathrm{Re}(Z_p)<0$ for all $p\leq\kappa$,
and consequently $B<0$.
\end{proposition}

\begin{proof}
Write $\mathrm{Re}(Z_p) = \mathrm{Re}(\Zt_p) + \mathrm{Re}(R_{p,N})$.
By~(A), $\mathrm{Re}(\Zt_p)<0$.
By~(B), $|\mathrm{Re}(R_{p,N})|<|\mathrm{Re}(\Zt_p)|$.
Hence $\mathrm{Re}(Z_p)<0$.
Then $B=\sum_{p\leq\kappa}(\log p)^2\mathrm{Re}(Z_p)<0$
by Proposition~\ref{prop:decomp}.
\end{proof}

Assumption~(A) is the Bias Conjecture (Open Problem~\ref{op:bias}).
Assumption~(B) requires quantitative truncation-error control via
the bound of Definition~\ref{def:trunc}, which is not established here.
Neither condition is verified in this paper.

The logical chain from Conjecture~\ref{conj:bias} to the Curvature Bias
at $\sigma=\tfrac{1}{2}$ proceeds in three steps, all conditional.
Under the Bias Conjecture and the truncation control of~(B),
Proposition~\ref{prop:sign} gives $B<0$.
By Reduction Lemma~\ref{lem:reduction}, this establishes the
negative curvature bias.
Complementing this with analytic control of the stationarity condition
$|O'(\tfrac{1}{2})|\ll W\cdot P$ (Open Problem~\ref{op:stat})
would then complete the $\sigma=\tfrac{1}{2}$ Curvature Bias argument,
subject to all stated auxiliary conditions.

% -------------------------------------------------------
\section{Open Problems}
\label{sec:open}

We list three open problems in decreasing order of centrality
to the program of this series.

\begin{openproblem}[Bias Conjecture]
\label{op:bias}
For each fixed $(\kappa,\varepsilon)$,
$\mathrm{Re}(\Zt_p(\varepsilon))<0$ for all primes $p\leq\kappa$.
This is the decisive open problem of this paper and of the series.
Its proof via quantitative estimates on the Guinand--Weil explicit
formula --- controlling $\Gamma$-terms and correction terms relative
to the candidate main term $-(\log p)/(2\sqrt{\pi}\,\varepsilon\sqrt{p})$
--- would, in conjunction with the truncation control of
Proposition~\ref{prop:sign}(B) and the Stationarity condition below,
complete the $\sigma=\tfrac{1}{2}$ Curvature Bias argument.
\end{openproblem}

\begin{openproblem}[Stationarity]
\label{op:stat}
Establish analytically that $|O'(\tfrac{1}{2})|\ll W\cdot P$,
where $W:=\sum_k e^{-\varepsilon^2\gamma_k^2}$ and $P=\pi(\kappa)$.
Numerically, $O'(\tfrac{1}{2})=-2.475$, approximately $11\%$ of $W\cdot P$
[Numerical].
An analytic proof of approximate stationarity of $O$ at $\sigma=\tfrac{1}{2}$
would complement the Bias Conjecture.
\end{openproblem}

\begin{openproblem}[Weighted spectral measure]
\label{op:measure}
Construct a weighted spectral measure $\nu^{(w)}_\kappa$ adapted to
$c_p$ and $\Estr$ such that
$\etaorig(\kappa) = 1 - m_1^{(w)}(\nu^{(w)}_\kappa)$
holds algebraically at finite $\kappa$.
The plain empirical measure $\nu_\kappa=(1/N)\sum_j\delta_{\mu_j}$ gives
$m_1(\nu_\kappa)=(1/N)\mathrm{Tr}(\Tt)$, which is a normalized trace
not identifiable with $c^T\Gun c/\Estr$.
The construction of an appropriate weighted version,
and its $\kappa\to\infty$ limit, remain open.
\end{openproblem}

% -------------------------------------------------------
\section{Non-Circularity}
\label{sec:nocirc}

We document explicitly that the results of this paper do not rely
on the objects they are directed toward.

\textbf{No RH as input.}
At no point is the Riemann Hypothesis assumed.
The ordinates $\gamma_k$ are used as numerically specified inputs.
No property of their distribution requiring RH is assumed or used.

\textbf{No GUE, Montgomery, or Hilbert--P\'olya hypothesis.}
The operator $\Tt(\sigma)$ is constructed explicitly from the coupling
map $\Phi(\sigma)$ (Definition~\ref{def:coupling}).
No random matrix hypothesis, no GUE pair-correlation conjecture,
and no Hilbert--P\'olya postulate enters the construction or the proofs.

\textbf{$\Tt(\sigma)$ is not postulated.}
It is derived as $\Phi(\sigma)\circ\Phi(\sigma)^*$ from an explicit matrix.
It is not a hypothetical object.

\textbf{$\Tt$ is not a Hilbert--P\'olya operator.}
As established in \S\,\ref{ssec:spectral}, the eigenvalues satisfy
$\mu_j\sim\CT/\gamma_{k(j)}\to 0$, which is incompatible with the
Hilbert--P\'olya requirement that eigenvalues grow like the zeros.

\textbf{Theorem~\ref{thm:trace} is purely algebraic.}
The trace formula is proved by a two-line trigonometric manipulation.
No asymptotic analysis and no analytic number theory are used.

\textbf{The Bias Conjecture is open.}
Conjecture~\ref{conj:bias} is clearly marked as an open problem.
No theorem in this paper implies it.

\textbf{$\gamma_k$ as inputs, not conclusions.}
The ordinates appear as numerically specified inputs to the model.
The claim $\mathrm{Re}(\rho_k)=\tfrac{1}{2}$ is not used anywhere;
the reference point $\sigma_0=\tfrac{1}{2}$ is chosen as the object
of study, not assumed to be the unique zero location.

% -------------------------------------------------------
\section*{Acknowledgments}

The author thanks the project's structured review and verification
workflow for systematic error detection across draft versions.

\emph{Pauca sed matura.} --- C.F.~Gauss

% -------------------------------------------------------
\section*{Call for Feedback}

This paper is part of the ongoing \textit{AnalysisLab\_L1--L5} project,
an open research initiative toward the Riemann Hypothesis.
All papers are freely available on Zenodo under CC~BY~4.0,
and the accompanying verification code is hosted on GitHub.

The author welcomes critical feedback, corrections, and collaboration:
Zenodo (\url{https://zenodo.org/search?q=Ulrich+Tehrani})
and GitHub (\url{https://github.com/utehrani/analysislab-nt}).

Topics of particular interest include: analytical approaches to the
Bias Conjecture via Guinand--Weil estimates; connections to existing
results on exponential sums over primes and zeros; and any errors or
gaps in the present arguments.

\noindent\textit{Pauca sed matura.} (C.F.\ Gauss)

% -------------------------------------------------------
\begin{thebibliography}{99}

\bibitem{Paper1}
U.~Tehrani,
``A Curvature Decomposition of the Explicit Formula'',
Zenodo, \url{https://doi.org/10.5281/zenodo.19025598},
April 2026, v8.

\bibitem{Paper2}
U.~Tehrani,
``From Local Curvature to the Weil Functional'',
Zenodo, \url{https://doi.org/10.5281/zenodo.19106992},
April 2026, v5.

\bibitem{Paper3}
U.~Tehrani,
``A Finite-Cutoff Hilbert-Space Model
for Prime--Zero Energy Structure'',
Zenodo, \url{https://doi.org/10.5281/zenodo.19307989},
April 2026, v2.

\bibitem{Paper4}
U.~Tehrani,
``A Dual Operator for Prime--Zero Coupling
and a Conditional Proof of Energy Asymmetry'',
Zenodo, \url{https://doi.org/10.5281/zenodo.19364703},
April 2026, v4.

\bibitem{Code}
U.~Tehrani,
\emph{AnalysisLab\_L1\_L5 --- Verification Scripts},
GitHub, \url{https://github.com/utehrani/analysislab-nt}.

\bibitem{Guinand}
A.~P.~Guinand,
``Fourier reciprocities and the Riemann zeta-function'',
\emph{Proc.\ London Math.\ Soc.}\ (2) \textbf{51} (1950), 401--414.

\bibitem{Weil}
A.~Weil,
``Sur les `formules explicites' de la th\'eorie des nombres premiers'',
\emph{Comm.\ S\'em.\ Math.\ Univ.\ Lund}, 1952, 252--265.

\bibitem{Mazhouda}
K.~Mazhouda,
``New asymptotic formulas for sums over zeros of functions
from the Selberg class'',
\emph{Ann.\ Sci.\ Math.\ Qu\'ebec}, 2019.
arXiv:1505.01544v2 [math.NT].

\end{thebibliography}

% -------------------------------------------------------
\bigskip\noindent\rule{\linewidth}{0.4pt}\smallskip
\noindent\small
\emph{Paper~5 v1.1 \textperiodcentered\ April~2026}\\
\emph{Pauca sed matura}

\end{document}
